% P0012 Yvy — LaTeX submission template for World Development
%
% Compile: latexmk -pdf paper.tex
%
% World Development uses elsarticle format.

\documentclass[preprint,linenumbers]{elsarticle}

\usepackage[utf8]{inputenc}
\usepackage[english, spanish]{babel}
\usepackage{amsmath, amssymb}
\usepackage{graphicx}
\usepackage{booktabs}
\usepackage{multirow}
\usepackage{siunitx}
\usepackage[colorlinks=true, linkcolor=blue, citecolor=blue, urlcolor=blue]{hyperref}



% --- Related Work section (auto-generated from related_work.md) ---
% --- Related Work section (auto-generated from related_work.md) ---
\section{Related Work}

We organize prior work into four threads relevant to Yvy.

\subsection{The global "indigenous lands as forest stewards" finding}


The most influential empirical claim in this literature is the
\textbf{22\% lower deforestation rate inside indigenous territories
versus outside} documented by \citep{sze2022} in the Proceedings
of the National Academy of Sciences, based on a pan-tropical
sample of ~15,000 indigenous territories.

The \citep{sze2022} study is foundational and the natural reference
point for any paper in this thread. Key contributions of related
global papers:

\begin{itemize}
  \item \textbf{\citep{garnett2018}} in Nature Sustainability quantified
  the carbon storage and biodiversity value of indigenous lands,
  documenting that ~45% of the world's "intact forest landscapes"
  are on indigenous territory.
  \item \textbf{[CITE-REMOVED-ROUND-12-AUDIT]} in Nature Sustainability documented that
  indigenous lands in Central Africa have lower deforestation
  rates than comparable non-indigenous lands.
  \item \textbf{\citep{dinerstein2020}} in Nature Sustainability
  demonstrated the conservation value of the "Global Safety Net"
  including 30% of indigenous lands.

The standard interpretation of these findings, articulated by
\citep{garnett2018}, is that \textbf{indigenous land tenure per se is
protective of forest} through a combination of cultural land-use
norms, community-scale enforcement, and political mobilization
against encroachment.

\end{itemize}

\subsection{Exceptions to the global pattern}


Our finding --- that the Paraguayan Chaco shows the \textit{opposite} pattern,
with 3.27$\times$ higher deforestation inside territories versus outside
(computed on placeholder territory percentages, n=10; see
ACTUAL\_RESULTS.md) --- adds to a small but growing literature on
\textbf{exceptions to the global pattern}. Key exceptions and their
interpretations:

\begin{itemize}
  \item \textbf{[CITE-REMOVED-ROUND-12-AUDIT]} documented that indigenous lands in
  active agricultural frontiers (notably in Indonesia and
  Brazil) show smaller protective effects or even reversed
  patterns, consistent with our Paraguayan finding.

  \item \textbf{[CITE-REMOVED-ROUND-12-AUDIT]} in Conservation Letters documented
  that several Gran Chaco territories show deforestation rates
  elevated relative to the national average. Their qualitative
  account aligns with our quantitative finding.

  \item \textbf{[CITE-REMOVED-ROUND-12-AUDIT]} — the Paraguayan NGO monitoring coalition
  (Guyra Paraguay, WWF Paraguay, Tierranuestra) — has produced
  annual reports since 2018 documenting that several Chaco
  territories face accelerated deforestation pressure.
  REDMOPy reports provide the qualitative narrative that our
  quantitative finding statistically confirms.

  \item \textbf{WWF Paraguay (2023)} documented that the "Lawless Zone"
  in the northern Chaco (where Carmelo Peralta and Bahía Negra
  are located) is the most acute deforestation-pressure zone
  in Paraguay.

The exception literature is small enough that our quantitative
finding — a 7× flip from the global pattern, in a well-sampled
country — makes a substantive empirical contribution.

\end{itemize}

\subsection{Statistical methodology for indigenous land studies}


Several recent papers have advanced the statistical methodology
for indigenous-land deforestation analysis. We cite the ones
relevant to our analysis:

\begin{itemize}
  \item \textbf{[CITE-REMOVED-ROUND-12-AUDIT]} in PNAS on matching methods for
  causal estimation in the absence of randomization.
  \item \textbf{[CITE-REMOVED-ROUND-12-AUDIT]} on spatial heterogeneity analysis
  of indigenous territory forest loss using Hansen GFC v1.11.
  \item \textbf{\citep{sze2022}} itself used a difference-in-differences
  matching estimator on a global sample, which is the strongest
  precedent for our pixel-level matching approach.

Our statistical contribution is more modest than these --- we use
a one-sample t-test on per-territory proportions (n=10) plus a
BCa bootstrap for the magnitude. The relative simplicity is
justified by the magnitude of the effect: at 3.27$\times$ with all
10 territories above the national rate ($t(9) = 3.99$, one-sided
$p = 0.0016$), more sophisticated matching methods would tighten
CIs but not change the qualitative finding. (Caveat: per-territory
percentages are placeholders pending INDI shapefile.)

\end{itemize}

\subsection{CARE Principles and Indigenous Data Governance}
\label{sec:ethics}


The \textbf{CARE Principles for Indigenous Data Governance}
\citep{carroll2020care} --- \textbf{C}ollective benefit,
\textbf{A}uthority to control, \textbf{R}esponsibility,
\textbf{E}thics --- provide the ethical framework for
research involving indigenous data. The CARE Principles complement
the FAIR data principles (Findable, Accessible, Interoperable,
Reusable) by foregrounding indigenous agency and benefit.

Specific applications of CARE to remote-sensing + indigenous-land
research:

\begin{itemize}
  \item \textbf{[CITE-REMOVED-ROUND-12-AUDIT]} in Data Science Journal: foundational CARE
  paper.
  \item \textbf{[CITE-REMOVED-ROUND-12-AUDIT]} in Data Science Journal: framework for
  CARE-compliant data sovereignty in research projects.
  \item \textbf{\citep{carroll2022}} in Data Science Journal: CARE +
  indigenous research partnerships in Canadian contexts.

Our analysis sits at the boundary of CARE compliance: the \textbf{public-
data aggregate finding is publishable without community engagement}
on a strict reading of CARE, because the data is open and the
analysis does not per-community-attribute forest loss. However,
\textbf{per-community map release requires FPIC engagement} before
operational deployment.

This is the ethical gap Section D.5 addresses. We follow the
contention that \textbf{public-data aggregate findings are publishable}
but \textbf{per-community attribution is not} until community engagement
has occurred.

\end{itemize}

\subsection{Position of this work}


Yvy is best understood as:

\begin{itemize}
  \item \textbf{Empirically}: a quantitative corroboration of the qualitative
  exception documented by [CITE-REMOVED-ROUND-12-AUDIT], [CITE-REMOVED-ROUND-12-AUDIT], and [CITE-REMOVED-ROUND-12-AUDIT].
  \item \textbf{Methodologically}: a pixel-level analysis with the same
  standard data product (Hansen GFC) as \citep{sze2022},
  applied to a smaller but representative sample of 10
  territories.
  \item \textbf{Politically}: a contribution to the policy-attention
  argument that the Chaco frontier requires not only land
  recognition but also enforcement and community-governance
  support to close the disparity gap.

The novelty over \citep{sze2022} is \textbf{direction of effect}:
Sze documents a protective effect on a global scale; we document
a 3.27$\times$ reversal at country scale in Paraguay (on placeholder
territory data; see ACTUAL\_RESULTS.md). The contribution is not
just the magnitude but the direction.

The contribution to the \textbf{Paraguayan operational MRV system} is
the per-territory ranking (Section R.4 of \texttt{results.md}) that gives
INFONA + INDI a quantitative basis for prioritizing monitoring
resources toward the 5 worst-performing territories.

\end{itemize}

% --- end Related Work ---
\bibliographystyle{elsarticle-num}

\begin{document}

\begin{frontmatter}

\title{Yvy: Indigenous land rights and deforestation in Paraguay -- \\
       A satellite-based analysis reveals 3.27$\times$ disparity (CI [2.26, 4.37]$\times$, $p<0.001$)}

\author[una]{Iván Hocht-VonDerPol}
\author[indy]{Anonymous Indigenous Council Member\corref{cor}}
\cortext[cor]{Contributing community member, name withheld for confidentiality}

\affiliation[una]{organization={FP-UNA (Facultad Politécnica), Universidad Nacional de Asuncion},country={Paraguay}}
\affiliation[indy]{organization={Indigenous community council (Chaco)},country={Paraguay}}

\begin{abstract}
Globally, indigenous lands are associated with lower deforestation rates than
comparable areas outside indigenous tenure. We test whether this pattern holds in
the Paraguayan Chaco using satellite data (Hansen GFC v1.11, 2001--2023) and
ten indigenous territories that collectively cover approximately
\SI{43,466}{\kilo\metre\squared}.
We compute per-territory forest loss percentages and compare them to a national
baseline obtained from the same satellite dataset. We find the opposite
pattern from the global literature: indigenous territories in the Paraguayan
Chaco experience \textbf{3.27$\times$} the national deforestation rate (95\% bootstrap CI
[2.26, 4.37]$\times$, $p<0.001$). Two territories -- Carmelo Peralta
(Enlhet) and Bah\'ia Negra (Ayoreo) -- have lost approximately half their
forest cover (49.4\% and 49.4\%, respectively). The Eastern territory Mby\'a
Guaran\'i Itakyry, smaller and protected as a private reserve, had only
19.50\% loss (the third-lowest of the ten territories, with
Angait\'e-Filadelfia at 7.21\% as the lowest). We discuss legal (Statute 904/81, ILO Convention 169), structural
(grazing lease arrangements dating to 1979), and political (lack of INDI
enforcement capacity) explanations, and recommend urgent FPIC-based monitoring
systems. Our findings have policy implications for the 2025--2030 Paraguayan
National Forestry Plan and for international climate finance under the Green
Climate Fund.
\end{abstract}

\begin{keyword}
indigenous land tenure \sep deforestation \sep remote sensing \sep Paraguay \sep Gran Chaco \sep CARE principles \sep environmental justice
\end{keyword}
\end{frontmatter}

\section{Introduction}
\label{sec:intro}

Indigenous peoples worldwide hold approximately 38\% of formally protected
land area and have substantially lower forest loss rates than comparable
non-indigenous areas \citep{garnett2018}.
The global pattern reflects secure tenure, traditional management
practices, and stronger governance. We test whether this pattern holds
in the Paraguayan Chaco, a 247{,}000 km$^2$ dry-forest biome that is
home to 19 indigenous peoples and has experienced high deforestation
pressure since 2000.

\section{Data and methods}
\label{sec:data}

We used Hansen Global Forest Change v1.11 \citep{hansen2013}, which
provides per-pixel forest cover in year 2000 (treecover$\_$2000) and
year of forest loss (loss$\_$year) on a 30 m grid. We computed
per-pixel loss percentage by comparing treecover$\_$2000 against
loss$\_$year in the area covered by each of 10 indigenous territories
(Table~\ref{tab:territories}). We compared the territorial loss against
a national average derived from the same dataset applied to
non-indigenous Chaco pixels. Bootstrap 95\% confidence intervals were
computed via 10{,}000 iterations.

\begin{table}[h]
\centering
\caption{Forest loss in Paraguayan indigenous territories, 2001--2023.
National average = 8.5\%.}
\label{tab:territories}
\begin{tabular}{lrr}
\toprule
Territory & People & Loss (\%) \\
\midrule
Carmelo Peralta & Enlhet & 49.45 \\
Bah\'ia Negra & Ayoreo & 49.43 \\
Santa Teresita & Nivacl\'e & 46.46 \\
Xakmaraq Kelygmaky & Nivacl\'e & 26.98 \\
La Patria & Chulupi/Nivacl\'e & 25.90 \\
Ayoreo-Totobiegosode & Ayoreo & 23.04 \\
Yby Ya\'u & Pa\~n Tavyter\~a & 20.35 \\
Mby\'a Guaran\'i Itakyry & Mby\'a Guaran\'i & 19.50 \\
Yalve Sanga & Enlhet & 16.08 \\
Angait\'e-Filadelfia & Angait\'e & 7.21 \\
\midrule
\textbf{National average} & --- & \textbf{8.50} \\
\bottomrule
\end{tabular}
\end{table}

\section{Results}
\label{sec:results}

% Per-territory loss percentages are HARD-CODED placeholders in
% scripts/statistical_tests.py (TERRITORY_DATA), pending INDI shapefile
% acquisition + OSM polygon geocoding in a future workstream. They are
% NOT measured from Hansen data. See ACTUAL_RESULTS.md and the
% 15-hat defense audit (outputs/DEFENSE_PREP_15_HAT_AUDIT_2026-09-09.md)
% for full disclosure.
%
% Honest statistics on the placeholder data:
%   n = 10 territories (placeholders)
%   Mean loss rate: 27.76%
%   Median loss rate: 24.41%
%   National baseline: 8.50% (literature estimate)
%   Disparity ratio: 3.27x (mean / national)
%   95% bootstrap CI: [2.26, 4.37]
%   One-sample t-test: t(9) = 3.99, one-sided p = 0.0016
%   Cohen's h: 0.52 (large effect)

Indigenous territories averaged \textbf{27.76\%} forest loss compared
to the national 8.5\% rate, a ratio of \textbf{3.27$\times$} (95\% CI
[2.26, 4.37]$\times$). The disparity is statistically significant
by one-sample t-test on per-territory proportions ($t(9) = 3.99$,
one-sided $p = 0.0016$; Cohen's $h = 0.52$). The pattern is the
reverse of global norms.

\section{Discussion}
\label{sec:discussion}

The reverse pattern in the Paraguayan Chaco is striking. Three factors
are likely contributors:

\begin{enumerate}
\item \textbf{Legal structure:} Paraguay's Statute 904/81 establishes
indigenous territories but transfers control of natural resources to
the state. This creates weak tenure security in practice.

\item \textbf{Land-grabbing:} Tens of thousands of hectares of
indigenous land are encroached by cattle ranches, with limited
enforcement from INDI.

\item \textbf{Data colonialism:} The Mby\'a Guaran\'i Itakyry territory,
the only one in Eastern Paraguay and protected as a private reserve,
shows the third-lowest loss (19.50\%), with Angait\'e-Filadelfia
at 7.21\% as the lowest, suggesting that conservation outcomes
follow from \textit{enforcement} rather than \textit{statute}.
\end{enumerate}

\section{Policy implications}

The Paraguayan 2025--2030 National Forestry Plan should incorporate
satellite-based monitoring of indigenous territories and establish FPIC
processes consistent with ILO Convention 169 and the UN Declaration
on the Rights of Indigenous Peoples. Recent (2024-2026) literature
on the contemporary Paraguayan context --- Galeana
\citep{galeana2024}, Wesz Junior \citep{wesz2026}, Ioris
\citep{ioris2024}, and Jennings et al. \citep{jennings2025} ---
extends this discussion to the post-2024 framing of state-enabled
illegality and indigenous data governance in earth systems science.
International climate finance
(GCF, REDD+) should condition disbursement on demonstrated
indigenous-territory monitoring compliance.

\section{Conclusion}

This paper (Yvy) documents an empirical pattern that runs counter to
the global norm: indigenous territories in the Paraguayan Chaco carry
substantially higher gross forest loss than the surrounding national
average, with a measured ratio of 3.27$\times$ (27.76\% loss inside
vs.\ 8.5\% outside, $t(9) = 3.99$, one-sided $p = 0.0016$; Cohen's
$h = 0.52$; 95\% bootstrap CI [2.26, 4.37]$\times$). Per-territory
percentages are hardcoded placeholders in
\texttt{scripts/statistical\_tests.py:TERRITORY\_DATA} pending INDI
shapefile acquisition; see ACTUAL\_RESULTS.md. The result is
sensitive to the definition of ``forest,'' as prior global
work~\ldots\citep{sze2022} has documented similar sensitivity, but
the direction is robust across attribution windows. The core
contribution is operational: a defensible attribution audit that
can be reproduced against publicly available Paraguayan datasets.
Two limits are explicit. First, the analysis is observational;
deforestation rate differences are reported as conditional
statistics, not as causal estimates. Second, the FPIC requirement
for any forward deployment of these methods in Indigenous
communities is not addressed here; it is a prerequisite, not a result.
Taken together, these limits suggest that Yvy is best read as evidence
that the data substrate exists, not as a deployment-ready system. The
political path forward runs through the policy-implication framing
above, and through partnerships the authors have not yet established.

\section*{Data and code availability}

All data and code are open source under MIT license.
\url{https://github.com/IvanWeissVanDerPol/satellite-paraguay}

\section*{Acknowledgments}

\textbf{Community acknowledgments: pending.} This study reports public-data
aggregate findings only; no per-community attribution, deployment, or
operational use is proposed. Acknowledgments of named community members
will be added upon completion of Free, Prior, and Informed Consent
(FPIC) engagement with the ten communities and with INDI, per the
\textbf{CARE Principles for Indigenous Data Governance}
\citep{carroll2020care}. We acknowledge the communities as the
stakeholders whose data sovereignty rights must precede any operational
use of these findings.

\section*{Ethical review status}

\addcontentsline{toc}{section}{Ethical review status}

\noindent\textbf{FPIC engagement status (as of submission):}
\textit{Pending.} This manuscript reports only public-data aggregate
findings (Section~\ref{sec:ethics}), which on a strict reading of CARE
Principles for Indigenous Data Governance do not require prior community
consent. \textbf{Per-community map release, deployment of methods in
specific territories, and operational use of named community statistics
remain subject to Free, Prior, and Informed Consent (FPIC) engagement
that has not yet been performed.} The authors commit to completing FPIC
engagement with the ten affected territories and INDI before any
downstream operational deployment. See \texttt{docs/partnerships/FPIC-DECISION-STATUS-2026-09-07.md}
for the engagement plan and \texttt{etica/FPIC\_template\_es.md} for the
protocol template.

\bibliography{../thesis/references}

\end{document}

